% v24. Four candidate signs are reduced to a perfect matching by the new
% integer and positive-cycle obstruction. Legacy labels are retained.
% No denominator cutoff is used.

\subsection{三高分母の完全マッチング定理と Schur 境界}
\label{sm14:section}

\begin{proposition}[三高分母に残る四つの正符号型]
\label{sm14:four-sign-types}
$q_0=1$，$q_1,q_2,q_3>1$ とし，整数 $p_{ij}$ が
\[
 \sum_{j\ne i}p_{ij}=-q_i,
 \qquad p_{ij}p_{ji}=q_iq_j-1
\]
を満たすとする．$p_{ij}>0$ のとき辺 $ij$ を置くグラフを $G^+$ とする．
このグラフは次の四型のいずれかである：
完全マッチング，端点に $0$ を持つ $P_3$ と孤立頂点の和，
内部頂点に $0$ を持つ $P_4$，または $0$ を含む三角形と孤立頂点の和．
標識付きグラフとしての個数はそれぞれ $3,6,6,3$ である．
この必要条件には，分母の奇数性，正の weights，接着，$D$ の相殺，node 条件を用いない．
\end{proposition}

\begin{proof}
pair products が正なので，全ての $p_{ij}$ は非零で，両方向は同符号である．
負の row sum から $G^+$ の最大次数は二以下である．
また，頂点 $0$ の三つの entries が全て負なら，その和は $-3$ 以下となるので，
$\deg_{G^+}(0)\geq1$ である．二つの孤立頂点 $i,j$ があるなら
\[
 |p_{ij}|\leq q_i-2,\qquad |p_{ji}|\leq q_j-2
\]
となり，$p_{ij}p_{ji}=q_iq_j-1$ に反する．

さらに $p_{0i}=-t<0$，$p_{i0}=-T<0$ なら
$q_i=tT+1>T$ である．従って頂点 $i$ の残る二つの entries がともに正になることはない．
これにより，$0$ が端点である正の $P_4$ と，正の $C_4$ は排除される．

残る除外は，$0$ が $P_3$ の中心である場合である．正の辺を $01,02$，孤立頂点を $3$ とし，
\[
 (p_{01},p_{02},p_{03})=(s,t,-u),\qquad
 (p_{10},p_{20},p_{30})=(S,T,-U)
\]
と書く．全ての文字は正整数で，$u=s+t+1$ である．
$q_1=sS+1$，$q_2=tT+1$，$q_3=uU+1$ となる．
$L=-p_{13}$，$N=-p_{23}$，$M=-p_{31}$，$O=-p_{32}$ と置くと，
二つの負の entries $p_{12},p_{21}$ と row sums から
\[
 L\leq(s+1)S,\quad N\leq(t+1)T,
 \quad M+O=(s+t)U+1
\]
を得る．一方，pair products は
$LM>suSU$，$NO>tuTU$ を与えるので，
\[
 M+O>uU\left(\frac{s}{s+1}+\frac{t}{t+1}\right)
 \geq(s+t+1)U\geq(s+t)U+1.
\]
中央の不等式は
\[
 u\left(\frac{s}{s+1}+\frac{t}{t+1}\right)-(s+t+1)
 =\frac{(s+t+1)(st-1)}{(s+1)(t+1)}\geq0
\]
による．これは row sum と矛盾する．
四頂点で最大次数二以下のグラフを上の条件と合わせると，記載した四型だけが残る．
\end{proof}

% v24. The four-ray isolated-vertex lemma shortens the v18 proof.
% The proof is unbounded; all integer-gap arguments are explicit.
\subsection{三高分母の正符号グラフは完全マッチングである}
\label{mh18:section}

命題~\ref{sm14:four-sign-types} の四つの候補から，整数性によって二つの道型を除き，
正の weights によって三角形型を除く．これにより，以前は仮定していた完全マッチング性が
全ての正の三高分母模型について成り立つ．

\begin{theorem}[三高分母の完全マッチング定理]
\label{mh18:matching}
$q_0=1$，$q_1,q_2,q_3>1$ とし，整数 $p_{ij}$ が行和と pair products
\[
 \sum_{j\ne i}p_{ij}=-q_i,\qquad p_{ij}p_{ji}=q_iq_j-1
\]
を満たすとする．$p_{ij}>0$ のとき無向辺 $ij$ を置くグラフを $G^+$ とする．
このとき次が成り立つ．
\begin{enumerate}
\item $G^+$ は完全マッチング，または低 ray $0$ を含む三角形と孤立頂点の和である．
\item 正の basin weights が存在することと，$G^+$ が完全マッチングであることは同値である．
\end{enumerate}
特に正の模型の低 row は，ちょうど二つの負の entry と一つの正の entry を持つ．
二つの負の low-row entries の相手同士が，もう一つの正の辺を作る．
分母の上限，整数 clutching，対向 $D$ の相殺，node 条件は仮定しない．
\end{theorem}

\begin{proof}
命題~\ref{sm14:four-sign-types} により，除くべき型は，$0$ を端点に持つ $P_3$ と孤立頂点，
$0$ を内部頂点に持つ $P_4$，および $0$ を含む三角形である．

\paragraph{$0$ が端点である $P_3$ の排除．}
正の辺を $03,13$ とし，正整数 $s,t,S,T,U$ により
\[
 (p_{01},p_{02},p_{03})=(-s,-t,s+t-1),\quad
 (p_{10},p_{20},p_{30})=(-S,-T,U)
\]
と書く．$u=s+t-1$ とすると，$q_1=sS+1,q_2=tT+1,q_3=uU+1$ である．
さらに
\[
 p_{13}=L>0,\quad p_{31}=M>0,\quad
 (p_{12},p_{21},p_{23},p_{32})=(-H,-K,-N,-O)
\]
と置き，
\[
 a=q_1q_2-1,\quad b=q_2q_3-1,\quad d=q_1q_3-1,
\]
\[
 B=q_1-S>0,\quad C=q_3+U>0,\quad k=q_2-T>0
\]
と略記する．行和と pair equations は
\[
 H=L+B,\quad O=M+C,\quad K+N=k,\quad
 HK=a,\ NO=b,\ LM=d
\]
を与える．従って
\begin{equation}\label{mh18:path-polynomial}
 (kC-b)L^2+\{kd-aC+B(kC-b)\}L+d(kB-a)=0.
\end{equation}
一方，直接計算により
\[
 kC-b=-U(sT-1)-T+1\le0,
\]
\[
 kB-a=-uST-S-T+1<0,\qquad
 kd-aC=-dT-aU+q_3-q_2<0.
\]
最後の不等式は $d=q_1q_3-1>q_3$ と $T\ge1$ による．
式~\eqref{mh18:path-polynomial} は $L>0$ で狭義負となるため矛盾する．

\paragraph{$0$ が内部頂点である $P_4$ の排除．}
正の辺を $01,02,13$ とし，今度は
\[
 (p_{01},p_{02},p_{03})=(s,t,-s-t-1),\quad
 (p_{10},p_{20},p_{30})=(S,T,-U)
\]
と書く．$u=s+t+1$，$q_1=sS+1,q_2=tT+1,q_3=uU+1$ とし，
$L,M,H,K,N,O,a,b,d$ は前段落と同じ符号で定める．
\[
 B=q_1+S,\quad C=q_3-U,\quad k=q_2+T
\]
と置けば，全て正であり，同じ式~\eqref{mh18:path-polynomial} が必要となる．
この場合は
\[
 kC-b=U(sT-1)+T+1>0,
\]
\[
 kB-a=uST+S+T+1>0,\qquad
 kd-aC=dT+aU+q_3-q_2>0.
\]
最後は $a=q_1q_2-1>q_2$，$U\ge1$ による．
式の左辺は $L>0$ で狭義正となり，再び矛盾する．
ここまでには正の weights を用いていない．

\paragraph{三角形に対する正値性の障害．}
三角形と孤立頂点の和は，補題~\ref{fg24:no-isolated} により正の weights を持たない．
この補題は，孤立頂点の負の row から作る二つの整数 gap の和が一になることを用いる．
各 row で非正の gap を選ぶと，gain が一以上の有向閉路が生じるためである．
従って低 ray を含む三角形型も排除される．

\paragraph{完全マッチングの十分性．}
正の二辺は互いに端点を共有しない．各正の辺 $ij$ について
\[
 \frac{p_{ij}}{q_i}<\frac{w_j}{w_i}<\frac{q_j}{p_{ji}}
\]
という開区間は，$p_{ij}p_{ji}=q_iq_j-1<q_iq_j$ により空でない．
二辺について独立に有理数の比を選べば，正の有理 weights が得られる．
負の entries の basin inequalities は自動的に成立する．これで同値性も従う．
\end{proof}

\begin{remark}[正値性の仮定と適用範囲]
\label{mh18:positivity-essential}
正値性を落とすと三角形型は実際に存在する．例えば
\[
 q=(1,5,19,31),\qquad
 P=\begin{pmatrix}
 0&1&1&-3\\
 4&0&2&-11\\
 18&47&0&-84\\
 -10&-14&-7&0
 \end{pmatrix}
\]
は全ての行和・pair products を満たす．しかし正の閉路 $0\to2\to1\to0$ の gain は
$188/95>1$ であり，正の weights は存在しない．
完全マッチング定理は，三つの高 $D$ が同時に非負でないことをまだ主張しない．
一方，\ref{sh14:matching-setup} の Schur 表示と
定理~\ref{sm14:beta-negative} の $\beta\le-1$ は，今後，全ての正の三高分母模型に適用できる．
\end{remark}


定理~\ref{mh18:matching} により，正の weights を持つ三高分母の整数模型は
全て完全マッチング型である．以下の Schur 表示は，したがって全ての正の三高分母模型を覆う．
行列の整数 row/pair equations と正値性の役割は区別され，
正値性を除いた三角形型は実際に存在する．

以下では完全マッチングに合わせて標識を選び，\ref{sh14:matching-setup} の記号
$s,t,S,T,U\geq1$，$u=s+t-1$，$k\geq0$，$\alpha,\beta\in\mathbb Z$ を用いる．
\[
 A=s\alpha-t,\qquad B=t\beta-s,\qquad C=s+uk
\]
と略記する．式~\eqref{sh14:matching-linear} は
\begin{equation}\label{sm14:linear}
 BS+AT=-\alpha\beta,\qquad
 -CS+(A+1)U=k(\alpha+1),\qquad
 (C-1)T+(B+1)U=-(\beta+1)(k+1)
\end{equation}
である．

\begin{lemma}[基本的な整数 bound]
\label{sm14:basic-bounds}
この表示では $\alpha\geq1$，$A\geq1$，$B\leq-2$，$C\geq2$ が成り立つ．
\end{lemma}

\begin{proof}
$\alpha+1\leq0$ なら，正の entries の絶対値 $L=uS+\alpha+1$，$M=sU-k$ は
$LM\leq suSU=(q_a-1)(q_c-1)<q_aq_c-1$ を満たして矛盾する．
従って $\alpha\geq0$ であり，第二式から $A\geq0$，さらに $s\alpha\geq t\geq1$ から
$\alpha\geq1$ を得る．$\beta\leq-1$ なら $B\leq-s-t\leq-2$，
$\beta\geq0$ なら第三式と $C\geq1$ から $B+1<0$ を得る．
$A=0$ なら第一式から $\beta>0$ だが，$t=s\alpha$ より
$B=s(\alpha\beta-1)\geq0$ となって矛盾する．
$C=1$ なら $s=1,k=0$ となり，第三式は $t\beta U=-\beta-1$ となる．
整数 $\beta\geq0$，$\beta\leq-1$ のどちらでも両辺の符号が矛盾する．
\end{proof}

\begin{theorem}[Schur 補助変数の符号]
\label{sm14:beta-negative}
完全マッチング型の正の整数模型では，canonical な向き $k\geq0$ において
\[
 \beta\leq-1
\]
が成り立つ．分母の上界，奇数性，接着および node 条件は不要である．
\end{theorem}

\begin{proof}
$\beta\geq0$ と仮定し，$b=-B=s-t\beta\geq2$ と置く．
$C=s+uk\geq s\geq b$ なので
\[
 \Delta=b(A+1)-C(A+b)\leq b(1-b)<0.
\]
三つの線形式から $S,U$ を消去すると，$\Delta T=N$，ただし
\[
 N=b(A+1)(\beta+1)(k+1)
   -(b-1)\bigl(C\alpha\beta+bk(\alpha+1)\bigr)
\]
となる．$A=s\alpha-t$，$b=s-t\beta$ を代入すると，恒等式
\begin{align*}
 N-\Delta
  &=(s+\beta)A+s+b(b+\beta-1)\\
  &\quad+k\bigl\{uA+\alpha[s+\beta(s+b-1)]+b(\beta+1)\bigr\}
\end{align*}
を得る．右辺は狭義正である．従って $T=N/\Delta<1$ となり，
$T$ が正整数であることに反する．
\end{proof}

\begin{theorem}[$\beta=-1$ の境界と高係数の負性]
\label{sm14:beta-boundary}
同じ完全マッチング型で $\beta=-1$ なら，必ず $k=0$ であり，$D_0>0$ となる．
従って三つの高分母係数の少なくとも一つは負である．
この結論にも分母上限，奇数性，接着および node 条件を用いない．
\end{theorem}

\begin{proof}
ここでは $b=s+t$，$u=b-1$，$B=-b$ と書く．$k\geq1$ と仮定すると
$C=s+uk\geq b$ であるから
\[
 F=C(A+b)-b(A+1)\geq b(b-1)>0.
\]
線形式~\eqref{sm14:linear} を解くと
\[
 T=\frac{u[(\alpha+b)k-s\alpha]}{F},\qquad
 S=\frac{u(A-\alpha)k-\alpha(s-1)(A+1)}{F}
\]
となる．$S>0$ から $A>\alpha$ が必要である．ところが
\[
 F-u[(\alpha+b)k-s\alpha]
 =(s-1)(s\alpha+b)+t^2+uk(A-\alpha)>0
\]
なので $T<1$ となり，整数性に反する．従って $k=0$ である．

このとき $C=s\geq2$ であり，第三式は $uU=(s-1)T$ となる．
$G=tA-b(s-1)$ と置き，第二式と第一式も合わせると
$GT=\alpha su>0$ を得る．従って $G>0$ であり，特異な場合 $G=0$ は存在しない．
その上で解を表示すると
\[
 T=\frac{\alpha su}{G},\quad
 U=\frac{\alpha s(s-1)}{G},\quad
 S=\frac{\alpha(s-1)(A+1)}{G}.
\]
$G>0$ から
$A+1>s+s(s-1)/t>s$ であり，$sS=(A+1)U$ により $S>U$ を得る．
従って $D_0=S+T-U-1>0$ である．全係数の和は零だから，高係数の和は負となる．
\end{proof}

\begin{corollary}[$\beta=-1$ の全整数模型を与える有限な約数手順]
\label{sm14:beta-boundary-factorization}
任意の固定した $s\geq2,t\geq1$ に対し，$\beta=-1$ の完全マッチング型は，
\[
 K_0=s^2+st+t^2-s-t,\qquad g=\gcd(s-1,t)
\]
を用いて次の有限手順で全て得られる．正の約数 $G\mid K_0g$ のうち
$\alpha=(G+K_0)/(st)$ が正整数になるものを取り，
\[
 S=\frac{\alpha(s-1)(s\alpha-t+1)}G,\qquad
 T=\frac{\alpha s(s+t-1)}G,\qquad
 U=\frac{\alpha s(s-1)}G
\]
が全て整数となるものだけを残す．$k=0$ として canonical 行列を復元する．
奇分母を要求する場合は，さらに $sS,tT,(s+t-1)U$ が全て偶数である条件を課す．
\end{corollary}

\begin{proof}
前定理の $G$ は $st\alpha-K_0$ である．その表示から
\[
 tT-(s+t-1)=\frac{K_0(s+t-1)}G,
 \qquad tU-(s-1)=\frac{K_0(s-1)}G
\]
となる．左辺は整数なので
$G\mid K_0\gcd(s+t-1,s-1)=K_0g$ である．
逆に，記載した検査を通ったデータは線形式を満たし，$S,T,U>0$ である．
canonical entries は
\[
 H=tS+\alpha,\quad K=sT-1,\quad
 L=uS+\alpha+1,\quad N=uT,\quad M=sU,\quad O=tU+1
\]
で全て正であり，row sums と pair equations が成り立つ．
完全マッチングの二辺に中点 weights を取れば正値性も従う．
よって候補生成にも棄却にも人工的な分母上限を置いていない．
\end{proof}

以上により，完全マッチング型で三つの高 $D$ が全て非負となる可能性を調べる際は
$\beta\leq-2$ に制限できる．完全マッチング性は定理~\ref{mh18:matching} で証明されたが，
残る $\beta\leq-2$ における三つの高 $D$ の同時非負性は未解決である．
